\chapter{Spectral theory in Banach algebras}

In the unital results of this chapter, $1$ abbreviates the identity $1_\A$ of
the ambient algebra $\A$; we write $1_\A$ when the space to which an expression
belongs needs to be specified.

\section{Banach algebras, involution, and basic examples}

\begin{cthdefinicion}{(Banach algebra)}
A \emph{Banach algebra} is an associative algebra $\A$ over $\mathbb C$, equipped with a
norm such that $(\A,\|\cdot\|)$ is a Banach space and
\[
  \|ab\|\le \|a\|\,\|b\|,\qquad a,b\in\A.
\]
We say that $\A$ is \emph{unital} if there is an identity $1_\A$ with $\|1_\A\|=1$, and \emph{commutative} if
$ab=ba$ for every $a,b\in\A$.
\end{cthdefinicion}
\origen{Definition 1.1}

\begin{cthejemplo}{(Bounded functions as a Banach algebra)}
We seek a commutative model of the definition. Let $X$ be nonempty,
locally compact, and Hausdorff, and consider
\[
 C_b(X)=\{f:X\to\mathbb C: f\text{ is continuous and bounded}\},\qquad
 \|f\|_\infty=\sup_{x\in X}|f(x)|,
\]
with pointwise operations. For $f,g\in C_b(X)$,
$|f(x)g(x)|\le\|f\|_\infty\|g\|_\infty$ at each $x$, and therefore
$\|fg\|_\infty\le\|f\|_\infty\|g\|_\infty$.
The uniform limit of continuous bounded functions remains in $C_b(X)$,
which gives completeness; the constant function $1$ is the identity.
Thus $C_b(X)$ meets the definition of a Banach algebra and is commutative.
If $X$ is compact, every continuous function is bounded and $C_b(X)=C(X)$.
\end{cthejemplo}

\begin{cthejemplo}{(Composition of bounded operators)}
To contrast with pointwise multiplication, let $X$ be a nonzero Banach space.
We study $\mathcal B(X)$ (also denoted $\mathcal L(X)$), the algebra
of bounded linear operators $X\to X$, with composition and the operator
norm. The inequality $\|STx\|\le\|S\|\|T\|\|x\|$ implies
$\|ST\|\le\|S\|\|T\|$; moreover, the identity $I_X$ is a unit and
the norm limit of bounded operators remains in $\mathcal B(X)$.
The product can change when the order is reversed: in
$\mathcal B(\mathbb C^2)=M_2(\mathbb C)$, if $E_{ij}$ denotes the matrix
with a $1$ in position $(i,j)$ and zeros elsewhere,
$E_{12}E_{21}=E_{11}$, whereas $E_{21}E_{12}=E_{22}$.
This example shows why subsequent proofs cannot assume
commutativity of the entire algebra.
\end{cthejemplo}

\begin{investigacion}
\textbf{Connection with the literature: a direct example.}
Under Karlovich's hypotheses for a starlike curve $\Gamma$ and $1<p<\infty$, the boundary
problem leads to an element of $\mathcal B(L^p(\Gamma))$:
\[
 T=V_\alpha P_+ + M_G P_-\in \mathcal B(L^p(\Gamma)),
\]
where $V_\alpha f=f\circ\alpha$, $M_Gf=Gf$, and $P_\pm=(I\pm S_\Gamma)/2$.
The formula shows why sums and compositions are studied in a single algebra. In this
geometry one does not assume $S_\Gamma^2=I$ or that $P_\pm$ are algebraic projections;
the reduction by Sokhotski--Plemelj is explained at the end of the notes
\cite[pp.~463--464, eqs.~(1.1)--(1.2)]{Karlovich2017}.
\textbf{What to recognize:} $\mathcal B(L^p(\Gamma))$ realizes the abstract algebra of
bounded operators; the formula for $T$ combines elements of that algebra.
\end{investigacion}

\begin{cthdefinicion}{(Involution)}
An \emph{involution} on a complex algebra $\A$ is a map $a\mapsto a^*$ such that
\[
(\lambda a+b)^*=\overline\lambda a^*+b^*,\qquad
(a^*)^*=a,\qquad
(ab)^*=b^*a^*.
\]
These identities are required for all $a,b\in\A$ and $\lambda\in\mathbb C$.
If, in addition, $\A$ is a Banach algebra and $\|a^*\|=\|a\|$ for
every $a\in\A$, we call $\A$ an \emph{involutive Banach algebra}
in the sense of these notes.
\end{cthdefinicion}
\origen{Definition 1.4}

\begin{cthejemplo}{(Pointwise conjugation as an involution)}
In the algebra $C_b(X)$ from the preceding example, we want to identify a
concrete involution. For $f\in C_b(X)$ define
$f^*(x)=\overline{f(x)}$. If $f,g\in C_b(X)$, then
$(fg)^*(x)=\overline{f(x)g(x)}=g^*(x)f^*(x)$;
also $(f^*)^*=f$ and $\|f^*\|_\infty=\|f\|_\infty$.
Pointwise conjugation thus satisfies the definition of an isometric
involution; the equality $g^*f^*=f^*g^*$ reflects the commutativity of this model.
\end{cthejemplo}

\section{Invertibility and the Neumann series}

\begin{cthdefinicion}{(Group of invertible elements)}
Let $\A$ be a unital Banach algebra, with identity $1_\A$. Its group of invertible elements is
\[
 \GL(\A)=\{a\in\A: \exists b\in\A,\ ab=ba=1_\A\}.
\]
For $a\in\GL(\A)$, that $b$ is unique and is written $a^{-1}$.
\end{cthdefinicion}
\origen{Definition 1.7}

\begin{cthproposicion}{(Neumann series)}
Let $\A$ be a unital Banach algebra, with identity $1_\A$, and let $a\in\A$.
If $\|1_\A-a\|<1$, then $a\in\GL(\A)$ and
\[
 a^{-1}=\sum_{n=0}^{\infty}(1_\A-a)^n.
\]
\end{cthproposicion}
\origen{Proposition 1.8}

\begin{proof}[Proof]
Write $q=1_\A-a$. The hypothesis gives $\|q\|<1$ and, by submultiplicativity,
$\|q^n\|\le\|q\|^n$ for every $n\ge0$. Hence
\[
 \sum_{n=0}^{\infty}\|q^n\|
 \le \sum_{n=0}^{\infty}\|q\|^n<\infty.
\]
Completeness of $\A$ ensures that the partial sums
$s_N=\sum_{n=0}^{N}q^n$ converge in norm to an element $b\in\A$.
Since $a=1_\A-q$, multiplication on the left telescopes:
\[
 as_N=(1_\A-q)\sum_{n=0}^{N}q^n
 =\sum_{n=0}^{N}q^n-\sum_{n=1}^{N+1}q^n
 =1_\A-q^{N+1}.
\]
Multiplication on the right must also be checked:
\[
 s_Na=\sum_{n=0}^{N}q^n(1_\A-q)
 =\sum_{n=0}^{N}q^n-\sum_{n=1}^{N+1}q^n
 =1_\A-q^{N+1}.
\]
Here we have used only that $q$ commutes with its own powers, without requiring
$\A$ to be commutative. Finally,
$\|q^{N+1}\|\le\|q\|^{N+1}\to0$; continuity of
multiplication lets us pass to the limit in both identities to obtain
$ab=ba=1_\A$. Thus $a\in\GL(\A)$ and $b=a^{-1}$.
\end{proof}

\proofdiagram
{$\A$ unital Banach, $a\in\A$; $\|1_\A-a\|<1$}
{Completeness of $\A$; $\|(1-a)^n\|\le\|1-a\|^n$}
{The geometric series converges and the product telescopes}
{$a^{-1}=\sum_{n\ge0}(1_\A-a)^n$}

\begin{cthejemplo}{(Neumann series for a perturbed matrix)}
We want to invert a matrix close to the identity. In
$M_2(\mathbb C)$ with the operator norm, let $A=I_2-E$ with
$\|E\|<1$. The Neumann proposition applies because
$\|I_2-A\|=\|E\|<1$, and gives
\[
 A^{-1}=I_2+E+E^2+\cdots.
\]
For example, if
$E=\begin{pmatrix}0&t\\0&0\end{pmatrix}$ with $|t|<1$,
then $E^2=0$ and the series terminates at
$A^{-1}=I_2+E$; indeed, $(I_2-E)(I_2+E)=I_2$.
The norm condition ensures invertibility even when the series does not terminate.
\end{cthejemplo}

\begin{cthejemplo}{(Stable invertibility near an operator)}
Let $X$ be a nonzero Banach space and suppose an invertible operator
$A_0\in\mathcal B(X)$ is known. We want to invert $A_0+E$ when
$E\in\mathcal B(X)$ is small relative to $A_0$, in the sense that
$\|A_0^{-1}E\|<1$. The factorization
$A_0+E=A_0(I_X+A_0^{-1}E)$ allows us to apply Neumann to the second factor:
\[
 (A_0+E)^{-1}=(I_X+A_0^{-1}E)^{-1}A_0^{-1}
 =\sum_{n=0}^{\infty}(-A_0^{-1}E)^n A_0^{-1}.
\]
The formula computes an exact inverse and shows that the relevant scale of
the perturbation is $\|A_0^{-1}E\|$. The Fredholm question will be formulated
later through invertibility of a class in a quotient.
\end{cthejemplo}

\begin{cthnota}{(Where the hypotheses are used in the Neumann series)}
The inequality $\|1-a\|<1$ makes the series of norms summable; completeness of $\A$ produces
its sum inside $\A$; the identity lets us write $a=1-(1-a)$ and telescope.
Commutativity of the entire algebra is not used.
\end{cthnota}

\begin{cthproposicion}{(Local stability of invertibility)}
Let $\A$ be a unital Banach algebra. The set $\GL(\A)$ is open.
More precisely, if $a\in\GL(\A)$ and $b\in\A$ satisfy
\[
 \|b-a\|<\|a^{-1}\|^{-1},
\]
then $b\in\GL(\A)$.
\end{cthproposicion}
\origen{Remark 1.9}

\begin{proof}[Proof]
Put $c=a^{-1}(a-b)$. The hypothesis and submultiplicativity give
\[
 \|c\|\le \|a^{-1}\|\,\|a-b\|<1,
 \qquad b=a(1_\A-c).
\]
The preceding proposition gives a two-sided inverse
$(1_\A-c)^{-1}=\sum_{n\ge0}c^n$ inside $\A$. The product of
the two invertible factors is invertible, with
\[
 b^{-1}=(1_\A-c)^{-1}a^{-1}
 =\sum_{n=0}^{\infty}\bigl[a^{-1}(a-b)\bigr]^n a^{-1}.
\]
Therefore every ball $B(a,\|a^{-1}\|^{-1})$ is contained in
$\GL(\A)$, which is open.
\end{proof}

\proofdiagram
{$\A$ unital, $a\in\GL(\A)$ and $b\in\A$ close to $a$}
{Neumann series}
{Factor $b=a\,[1-a^{-1}(a-b)]$}
{$\|a^{-1}(a-b)\|<1\Rightarrow b$ invertible}

\begin{cthejemplo}{(Stability radius around $2I_n$)}
We want to make the invertibility ball of the preceding proposition concrete.
In $M_n(\mathbb C)$ with the induced operator norm, take
$a=2I_n$; then $a^{-1}=\frac12I_n$ and
$\|a^{-1}\|^{-1}=2$. If a matrix $b$ satisfies
$\|b-2I_n\|<2$, the local stability condition holds and
$b$ is invertible. Here the number $2$ measures how much the
matrix can vary before this criterion ceases to guarantee an inverse.
\end{cthejemplo}

\begin{cthcorolario}{(Invertibility outside the disk of radius $\|a\|$)}
Let $\A$ be a complex unital Banach algebra, with identity $1_\A$,
$a\in\A$, and $\lambda\in\mathbb C$. If $\|a\|<|\lambda|$, then
$\lambda1_\A-a\in\GL(\A)$.
\end{cthcorolario}
\origen{Corollary 1.10}

\begin{proof}[Proof]
The inequality $\|a\|<|\lambda|$ implies $\lambda\ne0$. For
$c=\lambda^{-1}a$ we have $\|c\|=\|a\|/|\lambda|<1$;
by the Neumann proposition, $1_\A-c$ has the two-sided inverse
$\sum_{n\ge0}c^n$. Since
$\lambda1_\A-a=\lambda(1_\A-c)$ and the scalar
$\lambda1_\A$ is invertible, so is $\lambda1_\A-a$.
More precisely, its two-sided inverse is
\[
 (\lambda1_\A-a)^{-1}
 =\lambda^{-1}\sum_{n=0}^{\infty}(\lambda^{-1}a)^n.
\]
The series converges in norm because $\|\lambda^{-1}a\|<1$.
\end{proof}

\proofdiagram
{$a\in\A$, $\lambda\in\mathbb C$; $|\lambda|>\|a\|$}
{$\lambda\ne0$ and $\|\lambda^{-1}a\|<1$}
{Neumann: $1_\A-\lambda^{-1}a\in\GL(\A)$}
{$\lambda1_\A-a\in\GL(\A)$}

\begin{cthejemplo}{(Invertibility of $2I_2-A$ outside the norm disk)}
Let us check the corollary with a matrix whose inverse can be seen
directly. Let
$A=\operatorname{diag}(1,-1)\in M_2(\mathbb C)$ with the operator norm;
$\|A\|=1$. For $\lambda=2$ we have $|\lambda|>\|A\|$ and
\[
  (2I_2-A)^{-1}=\operatorname{diag}(1,1/3).
\]
Diagonal multiplication verifies that this is a two-sided inverse.
The computation anticipates that there are no spectral values outside
the closed disk of radius $\|a\|$, a reading formalized after defining the spectrum.
\end{cthejemplo}

\section{Spectrum and resolvent}

\begin{cthdefinicion}{(Spectrum)}
Let $\A$ be a complex unital Banach algebra, with identity $1_\A$,
and let $a\in\A$. The \emph{spectrum of $a$ in $\A$} is
\[
 \spec_\A(a)=\{\lambda\in\mathbb C:\lambda1_\A-a\notin\GL(\A)\}.
\]
When there is no ambiguity, we write $\spec(a)$.
\end{cthdefinicion}
\origen{Definition 1.11}

\begin{cthobservacion}{(Spectral reading of the preceding corollary)}
Let $\A$ be a complex unital Banach algebra, $a\in\A$, and
$\lambda\in\mathbb C$. With the notation just defined, if
$|\lambda|>\|a\|$, then $\lambda\notin\sigma_\A(a)$.
\end{cthobservacion}

\begin{cthejemplo}{(Spectrum as the image of a continuous function)}
To relate the definition of spectrum to a familiar object, let
$X$ be nonempty compact Hausdorff and $f\in C(X)$. We want to decide when
$\lambda1-f$ is invertible in $C(X)$. If
$f(x)\ne\lambda$ for every $x\in X$, the function
$x\mapsto(\lambda-f(x))^{-1}$ is continuous and bounded by compactness,
and is its pointwise inverse. If $f(x_0)=\lambda$, no product with a
function in $C(X)$ can equal $1$ at $x_0$. Thus
$\spec_{C(X)}(f)=f(X)$: the spectrum coincides with the values
actually attained by the function.
\end{cthejemplo}

\pagebreak
\begin{cthproposicion}{(Compactness of the spectrum)}
Let $\A$ be a complex unital Banach algebra and $a\in\A$.
Then $\spec_\A(a)$ is compact and
\[
 \spec_\A(a)\subset\{\lambda\in\mathbb C:|\lambda|\le\|a\|\}.
\]
\end{cthproposicion}
\origen{Proposition 1.12}

\begin{proof}[Proof]
If $|\lambda|>\|a\|$, the preceding corollary gives
$\lambda1_\A-a\in\GL(\A)$. By the definition of spectrum, no
such $\lambda$ belongs to $\spec_\A(a)$; the spectrum is contained in the
closed disk of radius $\|a\|$ and is therefore bounded.
Moreover, the map $F:\mathbb C\to\A$, $F(\lambda)=\lambda1_\A-a$,
is continuous, since
$\|F(\lambda)-F(\mu)\|=|\lambda-\mu|\|1_\A\|$.
The stability proposition shows that $\GL(\A)$ is open, so
$F^{-1}(\GL(\A))=\mathbb C\setminus\spec_\A(a)$ is open.
The spectrum is closed and bounded in $\mathbb C$; by Heine--Borel, it is compact.
\end{proof}

\proofdiagram
{$\A$ unital, $a\in\A$; definition of $\spec_\A(a)$}
{Large scalars are in the resolvent set; $\GL(\A)$ is open}
{The spectrum's complement is open, and everything outside the disk lies in the resolvent set}
{$\spec(a)$ closed and bounded, hence compact}

\begin{cthejemplo}{(Spectrum of a triangular matrix)}
Let us verify the spectral bound for a matrix that is not diagonal.
Take $A=\begin{pmatrix}2&1\\0&3\end{pmatrix}\in M_2(\mathbb C)$.
The matrix $\lambda I_2-A$ is triangular and its determinant is
$(\lambda-2)(\lambda-3)$; it is therefore invertible exactly when
$\lambda\notin\{2,3\}$, and $\spec(A)=\{2,3\}$.
The compactness proposition also ensures $|\lambda|\le\|A\|$
for both values: the off-diagonal entry can change the
norm, but not these spectral values.
\end{cthejemplo}

\begin{cthdefinicion}{(Resolvent set and resolvent function)}
Let $\A$ be a complex unital Banach algebra and $a\in\A$.
The \emph{resolvent set} is $\rho_\A(a)=\mathbb C\setminus\spec_\A(a)$.
The \emph{resolvent function} is
\[
 R_a:\rho_\A(a)\to\A,\qquad R_a(z)=(z1_\A-a)^{-1}.
\]
When the ambient algebra is clear, we write $\rho(a)$.
\end{cthdefinicion}
\origen{Definition 1.13}

\noindent The conclusion of the corollary on invertibility outside the disk
can now also be written as
$|\lambda|>\|a\|\Rightarrow\lambda\in\rho_\A(a)$, because
$\rho_\A(a)=\mathbb C\setminus\spec_\A(a)$.

\begin{cthejemplo}{(Resolvent of a diagonal matrix)}
We want to compute the resolvent function of
$A=\operatorname{diag}(1,2)\in M_2(\mathbb C)$.
For $z\notin\{1,2\}$, the matrix $zI_2-A$ has nonzero diagonal entries,
and its inverse is found by inverting them separately:
\[
 R_A(z)=\operatorname{diag}\bigl((z-1)^{-1},(z-2)^{-1}\bigr).
\]
The formula is undefined precisely at the spectral values
$1$ and $2$; the resolvent thus records where invertibility fails.
\end{cthejemplo}

\begin{cthejemplo}{(The resolvent in the operator algebra)}
Let us place the same question in the algebra $\mathcal B(X)$ of a nonzero
Banach space $X$. Its identity is $I_X$, so for
$A\in\mathcal B(X)$ the preceding definitions read
\[
 \spec_{\mathcal B(X)}(A)
 =\{\lambda\in\mathbb C:\lambda I_X-A\notin\GL(\mathcal B(X))\},
 \qquad R_A(\lambda)=(\lambda I_X-A)^{-1}.
\]
Invertibility here means the existence of a \emph{bounded} linear inverse
$X\to X$. If $X=\mathbb C^2$ and
$A=\operatorname{diag}(1,2)$, we recover the calculation from the
preceding example. The abstract notation allows the same definition
to be applied to operators on infinite-dimensional spaces.
\end{cthejemplo}

\begin{cthlema}{(Resolvent identity and holomorphy)}
Let $\A$ be a complex unital Banach algebra and $a\in\A$.
The function $R_a:\rho_\A(a)\to\A$ is holomorphic and
\[
 R_a'(z)=-(z1-a)^{-2}.
\]
Moreover, for $z,w\in\rho(a)$,
\[
 R_a(z)-R_a(w)=(w-z)R_a(w)R_a(z).
\]
\end{cthlema}
\origen{Lemma 1.14}

\begin{proof}[Proof]
Write $b_z=z1_\A-a$ and $b_w=w1_\A-a$. Both are invertible
for $z,w\in\rho_\A(a)$, and $b_z-b_w=(z-w)1_\A$.
The general identity for two invertible elements,
$b_z^{-1}-b_w^{-1}=b_z^{-1}(b_w-b_z)b_w^{-1}$,
is checked because the right-hand side is
$b_z^{-1}b_wb_w^{-1}-b_z^{-1}b_zb_w^{-1}
=b_z^{-1}-b_w^{-1}$. Moreover, $b_z$
and $b_w$ commute, so their inverses commute as well. Thus
\[
 R_a(z)-R_a(w)=(w-z)R_a(z)R_a(w)
                =(w-z)R_a(w)R_a(z).
\]
To justify holomorphy as well, fix $w\in\rho_\A(a)$.
If $|z-w|\|R_a(w)\|<1$, the factorization
$b_z=b_w[1_\A+(z-w)R_a(w)]$ and Neumann give
\[
 R_a(z)=\sum_{n=0}^{\infty}(-1)^n(z-w)^nR_a(w)^{n+1}.
\]
This series converges in norm locally around $w$, so $R_a$
is holomorphic there. Dividing the preceding identity by $z-w$ and letting
$z\to w$ in the series yields
$R_a'(w)=-R_a(w)^2=-(w1_\A-a)^{-2}$.
\end{proof}

\proofdiagram
{$\A$ unital Banach, $a\in\A$; $z,w\in\rho_\A(a)$}
{Identity for inverses and local Neumann series}
{Obtain the resolvent identity and differentiate the local series}
{$R_a'(w)=-R_a(w)^2$}

\begin{cthejemplo}{(Derivative of the scalar resolvent)}
To check the sign in the lemma's formula, take
$\A=\mathbb C$ and $a=\alpha\in\mathbb C$.
For $z\ne\alpha$, $R_a(z)=(z-\alpha)^{-1}$ and
\[
 \frac{R_a(z+h)-R_a(z)}{h}
 =-\frac{1}{(z+h-\alpha)(z-\alpha)}
 \longrightarrow -(z-\alpha)^{-2}.
\]
This is the ordinary derivative of a reciprocal function; the lemma extends
the calculation to inverses in any unital Banach algebra.
\end{cthejemplo}

\begin{cthteorema}{(Nonemptiness of the spectrum)}
Let $\A$ be a complex unital Banach algebra and $a\in\A$. Then
\[
 \spec(a)\neq\varnothing.
\]
\end{cthteorema}
\origen{Theorem 1.15}

\begin{proof}[Proof]
Suppose $\spec_\A(a)=\varnothing$. Then
$R_a:\mathbb C\to\A$ is entire by the preceding lemma. In particular,
$R_a(0)=(-a)^{-1}\ne0$, since it is the inverse of an element of the
nonzero algebra. Hahn--Banach gives
$\varphi\in\A^*$ with $\varphi(R_a(0))\ne0$.
For $|z|>2\|a\|$, the Neumann series gives
\[
 R_a(z)=z^{-1}(1_\A-z^{-1}a)^{-1},\qquad
 \|R_a(z)\|\le\frac{1}{|z|}
 \sum_{n=0}^{\infty}\left(\frac{\|a\|}{|z|}\right)^n
 \le\frac{2}{|z|}.
\]
The scalar function $f=\varphi\circ R_a$ is entire and tends to zero
as $|z|\to\infty$. It is bounded outside a disk by the estimate,
and inside it by continuity and compactness. Liouville's theorem forces
$f$ to be constant; its limit at infinity makes that constant zero.
This contradicts $f(0)=\varphi(R_a(0))\ne0$. Hence
$\spec_\A(a)\ne\varnothing$.
\end{proof}

\proofdiagrambranch
{$\A$ complex unital Banach, $a\in\A$; suppose $\spec(a)=\varnothing$}
{$R_a$ would be entire and $R_a(z)\to0$ as $|z|\to\infty$}
{Hahn--Banach + Liouville applied to $\varphi\circ R_a$}
{Choose $\varphi$ with $\varphi(R_a(0))\ne0$; Liouville gives $\varphi\circ R_a=0$}
{$\spec(a)$ must be nonempty}

\begin{cthejemplo}{(Spectrum without using a determinant)}
As a familiar case, consider
$A=\begin{pmatrix}1&1\\0&1\end{pmatrix}\in M_2(\mathbb C)$.
The determinant of $\lambda I_2-A$ is $(\lambda-1)^2$, so
$\spec(A)=\{1\}$ and the spectrum is nonempty.
The nonemptiness theorem gives the same general conclusion without
using determinants, which are unavailable for an arbitrary
infinite-dimensional operator.
\end{cthejemplo}

\begin{cthnota}{(Where the hypotheses of nonemptiness are used)}
The complex scalar field allows Liouville's theorem to be applied to $\varphi\circ R_a$; completeness
guarantees the Neumann series and $\A$-valued holomorphy; Hahn--Banach separates a
nonzero value of $R_a$. The identity appears in $\lambda1-a$. Commutativity is
not needed.
\end{cthnota}

\section{Characters, maximal ideals, and the Gelfand transform}

\begin{cthteorema}{(Gelfand--Mazur rigidity)}
If $\A$ is a complex unital Banach algebra and every nonzero element
is invertible, then
$\A=\mathbb C1_\A$.
\end{cthteorema}
\origen{Theorem 1.16}

\begin{proof}[Proof]
Fix $a\in\A$. Nonemptiness of the spectrum gives
$\lambda\in\spec_\A(a)$ for some $\lambda\in\mathbb C$.
By definition, $\lambda1_\A-a$ is not invertible. Since every
nonzero element of $\A$ is invertible by hypothesis, necessarily
$\lambda1_\A-a=0$, that is, $a=\lambda1_\A$.
Thus the unital homomorphism $\mathbb C\to\A$,
$\lambda\mapsto\lambda1_\A$, is surjective; it is injective because
$1_\A\ne0$. It identifies $\A$ with $\mathbb C1_\A$.
\end{proof}

\proofdiagram
{$\A$ complex unital; every $a\in\A\setminus\{0\}$ is invertible}
{Nonemptiness of $\spec(a)$}
{Choose $\lambda\in\spec(a)$; then $\lambda1-a$ is not invertible}
{$\lambda1-a=0$, so $a=\lambda1$}

\begin{cthejemplo}{(The scalars as a division algebra)}
Let us test the theorem's hypothesis in the simplest model:
$\A=\mathbb C$ with absolute value. Every $z\ne0$ has inverse
$z^{-1}=1/z$, and the spectrum of $\alpha\in\mathbb C$ is
$\{\alpha\}$ because $\lambda-\alpha$ ceases to be invertible only
when it is zero. Gelfand--Mazur says that any complex unital
Banach algebra in which every nonzero element is invertible
is, algebraically, this same scalar model.
\end{cthejemplo}

\begin{cthdefinicion}{(Character and Gelfand spectrum)}
A \emph{character} of a commutative unital Banach algebra $\A$ is a nonzero
multiplicative linear functional $\varphi:\A\to\mathbb C$. The set of characters is denoted
\[
 \Omega(\A)=\{\varphi:\A\to\mathbb C:\varphi\text{ is a character}\},
\]
as a set. After continuity of the characters has been proved, it will be given
the weak-$*$ topology inherited from $\A^*$.
\end{cthdefinicion}
\origen{Definition 1.17}

\begin{cthejemplo}{(Point evaluation is a character)}
Let $X$ be nonempty compact Hausdorff and fix $x\in X$.
We want to check that evaluation at $x$ meets the definition of a character
of $C(X)$. The functional $\operatorname{ev}_x(f)=f(x)$ is linear and
\[
 \operatorname{ev}_x(fg)=f(x)g(x)
 =\operatorname{ev}_x(f)\operatorname{ev}_x(g),\qquad
 \operatorname{ev}_x(1)=1.
\]
It is therefore multiplicative and nonzero. Each point thus gives a
character; Assignment 2 will prove that $C(X)$ has no others.
\end{cthejemplo}

\begin{cthproposicion}{(Characters have norm one)}
Let $\A$ be a complex commutative unital Banach algebra and
$\varphi\in\Omega(\A)$. Then $\varphi$ is continuous and
\[
 \|\varphi\|=\varphi(1_\A)=1.
\]
\end{cthproposicion}
\origen{Proposition 1.18}

\begin{proof}[Proof]
Since $\varphi\ne0$, there is $b\in\A$ with $\varphi(b)\ne0$.
The equality $\varphi(b)=\varphi(1_\A b)=\varphi(1_\A)\varphi(b)$
forces $\varphi(1_\A)=1$. Fix $a\in\A$.
If $|\varphi(a)|>\|a\|$, the Neumann corollary shows that
$u=\varphi(a)1_\A-a$ is invertible. But
\[
 \varphi(u)=\varphi(a)\varphi(1_\A)-\varphi(a)=0.
\]
If $v=u^{-1}$, multiplicativity would give
$1=\varphi(1_\A)=\varphi(uv)=\varphi(u)\varphi(v)=0$,
a contradiction. Hence $|\varphi(a)|\le\|a\|$ for every
$a\in\A$. This bound proves continuity of $\varphi$ and
$\|\varphi\|\le1$; since $\|1_\A\|=\varphi(1_\A)=1$,
we obtain $\|\varphi\|\ge1$ and finally $\|\varphi\|=1$.
\end{proof}

\proofdiagram
{$\A$ commutative unital, $\varphi\in\Omega(\A)$, $a\in\A$}
{$\varphi(1)=1$; large scalars are in the resolvent set}
{If $|\varphi(a)|>\|a\|$, then $\varphi(a)1-a$ would be invertible but its image under $\varphi$ is $0$}
{$|\varphi(a)|\le\|a\|$ and $\|\varphi\|=1$}

\begin{cthejemplo}{(Norm of the evaluation functional on $C(X)$)}
Let $X$ be nonempty compact Hausdorff, fix $x\in X$, and consider
$\operatorname{ev}_x:C(X)\to\mathbb C$, $\operatorname{ev}_x(f)=f(x)$.
We want to calculate its norm. For any $f\in C(X)$,
\[
 |\operatorname{ev}_x(f)|=|f(x)|\le\|f\|_\infty,
\]
so $\|\operatorname{ev}_x\|\le1$. The constant function $1$
has uniform norm one and $\operatorname{ev}_x(1)=1$, hence
$\|\operatorname{ev}_x\|\ge1$. We conclude that
$\|\operatorname{ev}_x\|=1$, in agreement with the general proposition
on characters: evaluation at a point is a normalized continuous
functional.
\end{cthejemplo}

\begin{cthdefinicion}{(Topology of the character space)}
Let $\A$ be a complex commutative unital Banach algebra.
Since every $\varphi\in\Omega(\A)$ is continuous, $\Omega(\A)\subset\A^*$.
Its \emph{Gelfand topology} is the weak-$*$ topology
restricted from $\A^*$: a net $\varphi_i$ converges to $\varphi$ if
$\varphi_i(a)\to\varphi(a)$ for every $a\in\A$.
With this topology, $\Omega(\A)$ is compact and Hausdorff.
\end{cthdefinicion}
\origen{Definition 1.19}

\begin{proof}[Justification of compactness]
By the preceding proposition, every character has norm one, so
$\Omega(\A)$ lies in the unit ball of $\A^*$. If a net of
characters $\varphi_i$ converges weak-$*$ to $\varphi$ in that ball,
then $\varphi(1_\A)=\lim_i\varphi_i(1_\A)=1$, so
$\varphi\ne0$. Pointwise limits preserve linearity and, for
$a,b\in\A$,
\[
 \varphi(ab)=\lim_i\varphi_i(ab)
 =\lim_i\bigl(\varphi_i(a)\varphi_i(b)\bigr)
 =\varphi(a)\varphi(b).
\]
Thus $\Omega(\A)$ is closed in the weak-$*$ compact ball by
Banach--Alaoglu. The weak-$*$ topology is Hausdorff because two
distinct functionals differ on some element of $\A$.
\end{proof}

\begin{cthdefinicion}{(Two-sided ideal and quotient algebra)}
Let $\A$ be a Banach algebra. A subspace $J\subset\A$ is a
\emph{two-sided ideal} if $aJ\subset J$ and
$Ja\subset J$ for every $a\in\A$. Its classes $a+J$ form the quotient algebra
$\A/J$, with $(a+J)(b+J)=ab+J$. If $J$ is closed, the norm
\[
 \|a+J\|_{\A/J}=\inf_{j\in J}\|a+j\|
\]
makes $\A/J$ a Banach algebra; the proof of completeness is developed in
Assignment 1, Problem 4. A maximal ideal is a proper ideal that is not contained in another
larger proper ideal.
\end{cthdefinicion}

\begin{cthejemplo}{(Classes and norm in a simple quotient)}[unbreakable]
In arithmetic, $\mathbb Z/n\mathbb Z$ teaches us to identify elements
that differ by an ideal. Consider the normed version with
$X=\{p,q\}$ discrete, $\A=C(X)\cong\mathbb C^2$, and
$J=\{f:f(p)=0\}$. Two functions represent the same class
exactly when they have the same value at $p$. Moreover,
\[
 \|f+J\|_{\A/J}
 =\inf_{h\in J}\max\{|f(p)|,|f(q)+h(q)|\}
 =|f(p)|,
\]
since we can choose $h(q)=-f(q)$. The quotient retains only the
$p$ coordinate; the quotient norm measures the size of the
information that remains.
\end{cthejemplo}

\begin{cthproposicion}{(Maximal ideals and scalar quotients)}
Let $J$ be a maximal ideal of a unital Banach algebra $\A$.
Then $J$ is closed. If $\A$ is commutative, also
\[
 \A/J\cong\mathbb C.
\]
\end{cthproposicion}
\origen{Proposition 1.20}

\begin{proof}[Proof]
The closure $\overline J$ is still a two-sided ideal: if
$j_n\in J$ converges to $j$ and $a\in\A$, continuity of
multiplication gives $aj_n\to aj$ and $j_na\to ja$, both in
$\overline J$. If $J$ were not closed, maximality would require
$\overline J=\A$. Then there would be $j\in J$ with
$\|1_\A-j\|<1$. Neumann would make $j$ invertible, and
$1_\A=j^{-1}j\in J$, contradicting that $J$ is proper.
Hence $J$ is closed and $\A/J$ is a nonzero unital Banach algebra.
Its identity is $1_\A+J$ and has norm one: the quotient norm is at most
$\|1_\A\|=1$, while a distance less than one from $1_\A$ to $J$
would, by Neumann, produce an invertible element of $J$ and force
$J=\A$.

Now suppose $\A$ is commutative. The ideals of $\A/J$
correspond, under the quotient map, to the ideals
of $\A$ that contain $J$; by maximality, $\A/J$ has no
ideals other than $\{0\}$ and itself. Let $a+J\ne0$.
The ideal $(\A/J)(a+J)$ is nonzero and thus equals all of $\A/J$.
There is $b+J$ with $(b+J)(a+J)=1_\A+J$. Commutativity also gives
$(a+J)(b+J)=1_\A+J$, so every nonzero element
of the quotient is invertible. Gelfand--Mazur yields
$\A/J\cong\mathbb C$ as complex unital algebras.
\end{proof}

\proofdiagrambranch
{$\A$ unital Banach, $J\subset\A$ maximal}
{In the second assertion, $\A$ commutative}
{The closure of $J$ is still an ideal; $\A/J$ has no nontrivial ideals}
{Every nonzero element of $\A/J$ is invertible; apply Gelfand--Mazur}
{$J$ closed and $\A/J\cong\mathbb C$}

\begin{cthejemplo}{(A maximal ideal defined by evaluation)}[unbreakable]
Let $X$ be nonempty compact Hausdorff and fix $x\in X$.
We want to describe the quotient that forgets functions vanishing at $x$.
The set $J_x=\{f\in C(X):f(x)=0\}$ is the kernel of the character
$\operatorname{ev}_x$, and the map
\[
 C(X)/J_x\longrightarrow\mathbb C,\qquad f+J_x\longmapsto f(x)
\]
is well defined, multiplicative, and surjective because it includes
constant functions. Its kernel is zero, so it is an
isomorphism. Since the quotient is a field, $J_x$ is maximal:
the value at a point is all the information that survives.
\end{cthejemplo}

\begin{investigacion}
\textbf{Direct construction example.}
Let $X$ be an infinite-dimensional complex Banach space.
An operator $K\in\mathcal B(X)$ is \emph{compact} if
$K$ maps the unit ball to a relatively compact set. The set
$\mathcal K(X)$ of such operators is a closed two-sided ideal of $\mathcal B(X)$:
composition with a bounded operator preserves compactness, and the norm limit
of compact operators is compact. The quotient
\[
 \mathcal Q(X):=\mathcal B(X)/\mathcal K(X)
\]
is called the \emph{Calkin algebra} \cite[p.~2]{Karlovich2025}.
If $X$ is a Hilbert space, this quotient is also a
$C^*$-algebra; for a general Banach space, only its Banach algebra structure is used here.
\textbf{What to recognize:} $\mathcal K(X)$ plays exactly the role of the
closed ideal $J$ in the construction $\A/J$; $\mathcal Q(X)$ is the concrete quotient.
\end{investigacion}

\begin{investigacion}
\textbf{Connection with the literature: use of a quotient.}
Let $X$ be an infinite-dimensional complex Banach space.
An operator $A\in\mathcal B(X)$ is \emph{Fredholm} if its kernel is finite-dimensional,
its range is closed, and the codimension of its range is finite. The equivalence from
operator theory
\[
 A\text{ Fredholm}\quad\Longleftrightarrow\quad
 A+\mathcal K(X)\in\GL(\mathcal Q(X))
\]
identifies an operator question with a question of invertibility in a
quotient \cite[p.~2]{Karlovich2025}. In $X=\ell^2$, let $P$ be the rank-one projection
onto $\mathbb Ce_1$ and $A=I_X-P$. Then $A$ is not invertible, but
$A+\mathcal K(X)=I_X+\mathcal K(X)$ is: the difference $P$ disappears
in the quotient. This is an instructional elaboration based on the cited criterion.
\textbf{What to recognize:} the quotient $\A/J$ from Problem 1.4 reappears with
$\A=\mathcal B(X)$ and $J=\mathcal K(X)$; the equivalence with the Fredholm property is
an additional result from operator theory.
\end{investigacion}

\pagebreak
\begin{cthproposicion}{(Characters and maximal ideals)}
Let $\A$ be a complex commutative unital Banach algebra.
The assignment
\[
 \Omega(\A)\longrightarrow\{\text{maximal ideals of }\A\},\qquad
 \varphi\longmapsto\ker\varphi,
\]
is a bijection.
\end{cthproposicion}
\origen{Proposition 1.21}

\begin{proof}[Proof]
If $\varphi\in\Omega(\A)$, then $\varphi(1_\A)=1$ and
$\varphi$ is surjective because $\varphi(\lambda1_\A)=\lambda$
for every $\lambda\in\mathbb C$. Its kernel is a proper ideal, and the
first isomorphism theorem gives
$\A/\ker\varphi\cong\mathbb C$; since $\mathbb C$ has no
nonzero proper ideals, $\ker\varphi$ is maximal.

Conversely, if $J$ is maximal, the preceding proposition allows us to
choose a complex unital isomorphism
$\theta:\A/J\to\mathbb C$. Let
$q:\A\to\A/J$ be the quotient map. The composition
$\varphi=\theta\circ q$ is linear, multiplicative, and unital;
it is a character with $\ker\varphi=J$.
To check uniqueness, if $\psi$ is another character with the same
kernel, then for each $a\in\A$,
$a-\varphi(a)1_\A\in J=\ker\psi$. Therefore
$\psi(a)=\varphi(a)\psi(1_\A)=\varphi(a)$. Thus the two assignments
are inverses of one another.
\end{proof}

\proofdiagram
{$\A$ commutative unital; $\varphi\in\Omega(\A)$ or $J\subset\A$ maximal}
{Quotients $\A/J\cong\mathbb C$}
{A character's kernel has codimension $1$; conversely, use the quotient map}
{$\varphi\leftrightarrow\ker\varphi$}

\begin{cthejemplo}{(Recovering two points from ideals)}
To see the bijection between characters and maximal ideals,
take $X=\{p,q\}$ and $C(X)\cong\mathbb C^2$.
Evaluations give the distinct ideals
$J_p=\{f:f(p)=0\}$ and $J_q=\{f:f(q)=0\}$.
If $\varphi$ is any character, the indicator functions
$e_p,e_q$ satisfy $e_p+e_q=1$ and $e_pe_q=0$;
their images are scalar idempotents and exactly one equals $1$.
Since $f=f(p)e_p+f(q)e_q$, it follows that $\varphi$ is
$\operatorname{ev}_p$ or $\operatorname{ev}_q$.
In this case, the two points of the space are recovered as the
two maximal ideals of the algebra.
\end{cthejemplo}

\begin{cthteorema}{(The spectrum is read through characters)}
Let $\A$ be a complex commutative unital Banach algebra.
Then, for each $a\in\A$,
\[
 \spec(a)=\{\varphi(a):\varphi\in\Omega(\A)\}.
\]
\end{cthteorema}
\origen{Theorem 1.22}

\begin{proof}[Proof]
If $\lambda=\varphi(a)$ for a character $\varphi$, then
$\varphi(\lambda1_\A-a)=\lambda-\varphi(a)=0$.
An invertible element cannot have zero image under a unital character;
therefore $\lambda1_\A-a$ is not invertible and
$\lambda\in\spec_\A(a)$.

Conversely, let $\lambda\in\spec_\A(a)$ and put
$b=\lambda1_\A-a$. Since $b$ is not invertible, the principal ideal
$\A b$ is proper: if $1_\A=cb$ for some $c\in\A$,
commutativity would also give $bc=1_\A$. By Zorn's lemma,
$\A b$ is contained in a maximal ideal $J$: indeed, the union
of a chain of proper ideals containing $\A b$ remains
proper, since if it contained $1_\A$, that element would already belong
to one of its members. That union is an upper bound for the chain.
The preceding proposition gives a character $\varphi$ with $J=\ker\varphi$.
Since $b\in J$, we obtain
$0=\varphi(b)=\lambda-\varphi(a)$, whence
$\lambda=\varphi(a)$. Both inclusions are proved.
\end{proof}

\proofdiagrambranch
{$a\in\A$, $\varphi\in\Omega(\A)$ and $\lambda=\varphi(a)$}
{$\lambda\in\spec_\A(a)$}
{Characters $\leftrightarrow$ maximal ideals}
{$\lambda1-a$ lies in a maximal ideal exactly when some character annihilates it}
{$\lambda\in\spec(a)\Longleftrightarrow\exists\varphi\in\Omega(\A):\lambda=\varphi(a)$}

\begin{cthejemplo}{(Spectrum of the square function on $[0,1]$)}
Apply the spectral theorem via characters to
$X=[0,1]$ and $f(x)=x^2\in C(X)$.
Each evaluation $\operatorname{ev}_x$ takes on $f$ the value $x^2$,
and the pointwise invertibility criterion already checked for $C(X)$
rules out other spectral values. Therefore
\[
 \spec_{C(X)}(f)
 =\{\operatorname{ev}_x(f):x\in[0,1]\}
 =\{x^2:x\in[0,1]\}=[0,1].
\]
The spectrum recovers exactly the geometric range of the function.
\end{cthejemplo}

\begin{cthproposicion}{(Jacobson's lemma for $1-ab$)}
Let $\A$ be a unital Banach algebra with identity $1_\A$ and
$a,b\in\A$. Then
\[
 1_\A-ab\in\GL(\A)\quad\Longleftrightarrow\quad 1_\A-ba\in\GL(\A).
\]
If $1_\A-ab$ is invertible,
\[
 (1_\A-ba)^{-1}=1_\A+b(1_\A-ab)^{-1}a.
\]
\end{cthproposicion}
\origen{Proposition 1.23}

\begin{proof}[Proof]
Suppose $1_\A-ab$ is invertible and write
$c=(1_\A-ab)^{-1}$. The two inverse identities give
$c-abc=1_\A$ and $c-cab=1_\A$.
Define $d=1_\A+bca$. On the left,
\[
 (1_\A-ba)d
 =1_\A-ba+bca-babca
 =1_\A+b(c-1_\A-abc)a=1_\A.
\]
On the right, use the other identity for $c$:
\[
 d(1_\A-ba)
 =1_\A-ba+bca-bcaba
 =1_\A+b(c-1_\A-cab)a=1_\A.
\]
Thus $d$ is the two-sided inverse of $1_\A-ba$, proving
the formula in the statement. Starting instead from invertibility of
$1_\A-ba$, the same reasoning with $a$ and $b$ interchanged
proves the reverse implication.
\end{proof}

\proofdiagram
{$\A$ unital, $a,b\in\A$ and $1_\A-ab$ invertible}
{Associativity}
{Explicitly construct $1+b(1-ab)^{-1}a$}
{It is a two-sided inverse of $1-ba$}

\begin{cthejemplo}{(Jacobson's identity for noncommuting matrices)}[unbreakable]
We want to see the proposition without assuming the factors commute.
In $M_2(\mathbb C)$ take $A=E_{12}$ and
$B=\frac12E_{21}$, where $E_{ij}$ are the matrices with a
single entry $1$ in position $(i,j)$. Then
$AB=\frac12E_{11}$ and $BA=\frac12E_{22}$ are distinct, but
\[
 (I_2-AB)^{-1}=\operatorname{diag}(2,1),\qquad
 (I_2-BA)^{-1}=\operatorname{diag}(1,2).
\]
Jacobson's formula gives the second inverse from the first:
$I_2+B(I_2-AB)^{-1}A=I_2+E_{22}$.
Simultaneous invertibility follows from the structure of the product,
without requiring $AB=BA$.
\end{cthejemplo}

\begin{cthcorolario}{(Nonzero spectra of $ab$ and $ba$)}
Let $\A$ be a complex unital Banach algebra and $a,b\in\A$.
Then
\[
 \spec_\A(ab)\cup\{0\}=\spec_\A(ba)\cup\{0\},
 \qquad\text{equivalently}\qquad
 \spec_\A(ab)\setminus\{0\}=\spec_\A(ba)\setminus\{0\}.
\]
\end{cthcorolario}
\origen{Corollary 1.24}

\begin{proof}[Proof]
Let $\lambda\in\mathbb C\setminus\{0\}$ and use
the factorization
$\lambda1_\A-ab=\lambda(1_\A-\lambda^{-1}ab)$.
The scalar factor $\lambda1_\A$ is invertible, so
\[
\lambda1-ab\in\GL(\A)
\Longleftrightarrow
1-\lambda^{-1}ab\in\GL(\A)
\Longleftrightarrow
1-\lambda^{-1}ba\in\GL(\A)
\Longleftrightarrow
\lambda1-ba\in\GL(\A).
\]
The middle equivalence is the preceding lemma applied to
$\lambda^{-1}a$ and $b$; the last again uses the scalar factor
$\lambda\ne0$. By the definition of spectrum, $\lambda$ belongs
to $\spec_\A(ab)$ exactly when it belongs to
$\spec_\A(ba)$. Equality of the spectra away from zero
follows for every $\lambda\ne0$; adding $\{0\}$ to both
sets gives the first equality in the statement.
\end{proof}

\proofdiagram
{$a,b\in\A$, $\lambda\in\mathbb C\setminus\{0\}$}
{Jacobson's lemma}
{$\lambda1_\A-ab=\lambda(1_\A-\lambda^{-1}ab)$}
{The nonzero spectra coincide; adding $\{0\}$ gives the equality}

\begin{cthejemplo}{(Reversing the order across different dimensions)}
Let $A:\mathbb C^m\to\mathbb C^n$ and
$B:\mathbb C^n\to\mathbb C^m$. We want to compare the products
$AB\in M_n(\mathbb C)$ and $BA\in M_m(\mathbb C)$.
Jacobson's identity adapts to the two identities and, for
$\lambda\ne0$, gives
\[
 (I_m-\lambda^{-1}BA)^{-1}
 =I_m+\lambda^{-1}B(I_n-\lambda^{-1}AB)^{-1}A
\]
 if $I_n-\lambda^{-1}AB$ is invertible. As a concrete case, take
$A(t)=(t,0)$ from $\mathbb C$ to $\mathbb C^2$ and
$B(x,y)=x$ from $\mathbb C^2$ to $\mathbb C$. Then
$AB=\operatorname{diag}(1,0)$ and $BA=(1)$:
their nonzero spectra are both $\{1\}$, although only the first
contains $0$. This explains why the corollary adds $\{0\}$.
\end{cthejemplo}

\begin{cthdefinicion}{(Spectral radius)}
Let $\A$ be a complex unital Banach algebra and $a\in\A$.
The \emph{spectral radius} of $a$ is
\[
 r(a)=\max\{|\lambda|:\lambda\in\spec_\A(a)\}.
\]
\end{cthdefinicion}
\origen{Definition 1.25}

\begin{cthejemplo}{(Spectral radius and matrix norm)}
For $A=\operatorname{diag}(2,-3,i)\in M_3(\mathbb C)$,
the spectral values are its diagonal entries, so
$r(A)=\max\{2,3,1\}=3$. With the Euclidean operator norm, also
$\|A\|=3$, since it is diagonal. Equality is not automatic for
nonnormal matrices: if
$B=\begin{pmatrix}2&10\\0&-3\end{pmatrix}$,
its spectrum is still $\{2,-3\}$ and $r(B)=3$, whereas
$\|B\|\ge\|Be_2\|=\sqrt{109}>3$.
The radius measures spectral values and can be less than the norm.
\end{cthejemplo}

\begin{cthteorema}{(Spectral radius formula)}
Let $\A$ be a complex unital Banach algebra and $a\in\A$.
Then
\[
 r(a)=\lim_{n\to\infty}\|a^n\|^{1/n}
     =\inf_{n\ge1}\|a^n\|^{1/n}.
\]
\end{cthteorema}
\origen{Theorem 1.26}

\begin{proof}[Proof]
First check that the limit of norms exists. Put
$\beta_n=\|a^n\|$. Submultiplicativity gives
$\beta_{m+n}\le\beta_m\beta_n$. If $\beta_m=0$ for some
$m\ge1$, then $a^n=0$ for every $n\ge m$ and
$\lim_n\beta_n^{1/n}=\inf_{n\ge1}\beta_n^{1/n}=0$.
Otherwise, fix $m\ge1$ and write $n=qm+s$,
$0\le s<m$. Since
$\beta_n\le\beta_m^q\beta_s$ and there are only finitely many
possible $s$, taking $n$th roots and letting $n\to\infty$
gives
\[
 \limsup_{n\to\infty}\beta_n^{1/n}\le\beta_m^{1/m}.
\]
This holds for every $m$. On the other hand, each term
$\beta_n^{1/n}$ is at least $\inf_{m\ge1}\beta_m^{1/m}$.
Consequently,
\[
 L:=\lim_{n\to\infty}\|a^n\|^{1/n}
   =\inf_{n\ge1}\|a^n\|^{1/n}.
\]

Let us show that $r(a)\le L$. If $\lambda\in\spec_\A(a)$ and
$n\ge1$, the polynomial factorization
\[
 \lambda^n1_\A-a^n
 =(\lambda1_\A-a)
   \sum_{k=0}^{n-1}\lambda^{n-1-k}a^k
 \]
has commuting factors. If the product were invertible, each
factor would be: indeed, if $x$ and $y$ commute and $xy$ has
inverse $c$, the equality $x(xy)=(xy)x$ implies $xc=cx$ after
multiplication by $c$. Then
$x(yc)=(xy)c=1_\A$ and
$(yc)x=y(cx)=yx c=(xy)c=1_\A$, so $yc$ is
a two-sided inverse of $x$.
This contradicts the noninvertibility of $\lambda1_\A-a$. Thus
$\lambda^n\in\spec_\A(a^n)$, and the spectral bound already proved gives
$|\lambda|^n\le\|a^n\|$. Taking the maximum over the spectrum
and then taking roots yields $r(a)\le\|a^n\|^{1/n}$ for every
$n$, and therefore $r(a)\le L$.

For the other inequality, let $R=1/r(a)$ if $r(a)>0$ and
$R=\infty$ if $r(a)=0$. The function
$g(z)=(1_\A-za)^{-1}$ is defined and holomorphic on
$|z|<R$: for $z\ne0$ in that disk,
$|z^{-1}|>r(a)$ and
$1_\A-za=z(z^{-1}1_\A-a)$ is invertible; at $z=0$
it is the identity. If $z\ne0$, moreover,
$g(z)=z^{-1}R_a(z^{-1})$, which is holomorphic by the resolvent
lemma and holomorphy of $z\mapsto z^{-1}$ away from zero.
Near zero, Neumann gives
$g(z)=\sum_{n\ge0}a^nz^n$.
This series converges in norm for $|z|\|a\|<1$ and proves
holomorphy at zero as well.

Fix a real $t$ with $0<t<R$. For each $\varphi\in\A^*$,
$\varphi\circ g$ is holomorphic on $|z|<R$, and its Taylor series
at zero has coefficients $\varphi(a^n)$. Choose $s$ with
$t<s<R$. Cauchy's bound on the circle $|z|=s$ gives
$|\varphi(a^n)|\le M_{\varphi,s}s^{-n}$, where
$M_{\varphi,s}=\max_{|z|=s}|\varphi(g(z))|<\infty$.
Thus $t^n|\varphi(a^n)|\le M_{\varphi,s}(t/s)^n$ is
bounded for each $\varphi$. Apply Banach--Steinhaus on the Banach space
$\A^*$ to the continuous operators
$T_n:\A^*\to\mathbb C$,
$T_n(\varphi)=t^n\varphi(a^n)$. Pointwise boundedness implies
$\sup_n\|T_n\|<\infty$. Hahn--Banach gives
$\|T_n\|=t^n\|a^n\|$, so there is $C_t<\infty$ with
$t^n\|a^n\|\le C_t$ for every $n$. Taking roots and limits,
$L\le t^{-1}$. If $r(a)>0$, let $t\uparrow1/r(a)$ and
obtain $L\le r(a)$; if $r(a)=0$, we may take $t$
arbitrarily large and obtain $L=0$. Together with the first
inequality and the existence of $L$, this proves both equalities.
\end{proof}

\proofdiagrambranch
{$\A$ unital, $a\in\A$; $\lambda\in\spec(a)\Rightarrow\lambda^n\in\spec(a^n)$}
{Study $g(z)=(1-za)^{-1}$ on $|z|<R$, $R=1/r(a)$ or $\infty$}
{Holomorphy + Taylor + Banach--Steinhaus}
{Bound $t^n\|a^n\|$ for each $t<R$ and let $t\uparrow R$}
{$r(a)=\lim\|a^n\|^{1/n}$}

\begin{cthejemplo}{(The Volterra operator is quasinilpotent)}[unbreakable]
Let us study a nonzero operator whose spectral radius is zero.
On $C[0,1]$ with the uniform norm, define
$(Vf)(x)=\int_0^x f(t)\,dt$. Iterating the integral gives, for
$n\ge1$,
\[
 (V^nf)(x)=\frac1{(n-1)!}\int_0^x(x-t)^{n-1}f(t)\,dt,
 \qquad \|V^n\|\le\frac1{n!}.
\]
The bound follows by replacing $|f(t)|$ with $\|f\|_\infty$
and integrating the power.
By the spectral radius formula,
$r(V)=\lim_n\|V^n\|^{1/n}=0$.
Since the spectrum is nonempty and all its points have modulus
at most zero, $\spec(V)=\{0\}$.
This case distinguishes a quasinilpotent operator from the zero operator:
$V\ne0$ although its spectrum is $\{0\}$.
\end{cthejemplo}

\begin{cthnota}{(Interpretation of the formula for iterates)}
Let $X$ be a Banach space and $A\in\mathcal B(X)$. Then $\|A^n\|$
measures the largest effect of $n$ successive
applications. The limit $\lim_n\|A^n\|^{1/n}$ measures its asymptotic exponential rate and
equals $r(A)$. The Volterra example shows that this rate may be zero
even when $A\ne0$.
\end{cthnota}

\begin{cthdefinicion}{(Gelfand transform)}
Let $\A$ be a complex commutative unital Banach algebra.
The \emph{Gelfand transform} is
\[
 \Gamma:\A\to C(\Omega(\A)),\qquad
 a\mapsto\widehat a,\qquad \widehat a(\varphi)=\varphi(a).
\]
\end{cthdefinicion}
\origen{Definition 1.27}

\begin{cthejemplo}{(The transform recovers a function in $C(X)$)}
Take $X=[0,1]$ and $f(t)=t\in C(X)$ to make the definition concrete.
Each point $x\in X$ gives the evaluation character
$\operatorname{ev}_x$, and
\[
 \widehat f(\operatorname{ev}_x)
 =\operatorname{ev}_x(f)=x=f(x).
\]
Under the identification $\Omega(C(X))\cong X$, proved in Assignment 2,
the Gelfand transform is the function $f$ itself.
The formula shows how a function is read through all the
characters of the algebra.
\end{cthejemplo}

\begin{cthproposicion}{(The Gelfand transform is contractive)}
Let $\A$ be a complex commutative unital Banach algebra and
$a\in\A$. The Gelfand transform is linear, multiplicative,
unital, and
\[
 \|\widehat a\|_\infty=r(a)\le\|a\|.
\]
\end{cthproposicion}
\origen{Remark 1.28}

\begin{proof}[Proof]
For each $\varphi\in\Omega(\A)$, linearity and
multiplicativity of the character give
\[
 \widehat{\alpha a+b}(\varphi)
   =\alpha\widehat a(\varphi)+\widehat b(\varphi),\qquad
 \widehat{ab}(\varphi)=\widehat a(\varphi)\widehat b(\varphi),
 \qquad \widehat{1_\A}(\varphi)=1.
\]
Each $\widehat a$ is continuous in the weak-$*$ topology by its
definition. The theorem describing the spectrum through
characters identifies the range of $\widehat a$
with $\spec_\A(a)$. Hence
\[
 \|\widehat a\|_\infty=\sup_{\varphi\in\Omega(\A)}|\varphi(a)|
 =\max_{\lambda\in\spec(a)}|\lambda|=r(a)\le\|a\|.
\]
The maximum exists because the spectrum is compact and nonempty;
the last inequality is the spectral bound proved earlier.
\end{proof}

\proofdiagram
{Definition $\widehat a(\varphi)=\varphi(a)$}
{Linear and multiplicative characters; spectrum via characters}
{$\|\widehat a\|_\infty=\sup_{\varphi}|\varphi(a)|$}
{$\|\widehat a\|_\infty=r(a)\le\|a\|$}

\begin{investigacion}[unbreakable]
\textbf{Connection with the literature: a direct example of a commutative algebra.}
Karlovich considers $\SO(\mathbb R_+)$, the algebra of bounded continuous functions whose
oscillation on $[r,2r]$ tends to zero as $r\to0$ and as $r\to\infty$
\cite[pp.~465--466, §2]{Karlovich2017}. It is a commutative $C^*$-subalgebra of
$C_b(\mathbb R_+)$: it contains constants, is closed under addition, multiplication, and conjugation, and
the oscillation condition is preserved under uniform limits. The
$C^*$ identity will be recognized formally in Section 2; at this point all that is
needed is that the algebra is commutative and unital. For $a$ in this algebra,
$\widehat a(\xi)=\xi(a)$ evaluates $a$ at a character $\xi\in\Omega(\SO(\mathbb R_+))$.
Besides evaluations at $r>0$, the maximal ideal space has fibers over $0$ and
$\infty$. These fibers organize limiting behaviors that a single value $a(0)$ or
$a(\infty)$ would not describe. The Gelfand transform of this coefficient algebra
is a conceptual step toward Fredholm symbols, not the full Fredholm symbol.
\textbf{What to recognize:} the characters of $\A$ studied in the course
become ordinary evaluations and characters in the fibers of
$\Omega(\SO(\mathbb R_+))$ here. For the precise construction, see
\cite[§2, Propositions~2.1--2.2]{Karlovich2017}.
\end{investigacion}

\begin{cthdefinicion}{(Modular ideal)}
An ideal $I$ of an algebra $\A$ is called \emph{modular} if there is $u\in\A$ such that
\[
 a-ua\in I,\qquad a-au\in I,\qquad a\in\A.
\]
\end{cthdefinicion}
\origen{Definition 1.29}

\begin{cthejemplo}{(Local identity for an evaluation ideal)}
Let $X$ be locally compact Hausdorff and fix $x\in X$.
We want to check the modularity of the ideal
\[
 I_x=\{f\in C_0(X):f(x)=0\}.
\]
Choose a continuous compactly supported function $u$ with
$u(x)=1$, available by local compactness and the Hausdorff separation property;
in particular, $u\in C_0(X)$.
For every $f\in C_0(X)$ we have
$(f-uf)(x)=(f-fu)(x)=0$; therefore both differences belong to
$I_x$. The class $u+I_x$ acts as the identity of the quotient, although
$C_0(X)$ may lack a global identity.
\end{cthejemplo}

\begin{cthproposicion}{(Modularity criterion via the quotient)}
Let $\A$ be an algebra and $I\subsetneq\A$ a proper two-sided ideal.
Then $I$ is modular if and only if $\A/I$ has an identity.
\end{cthproposicion}
\origen{Proposition 1.31}

\begin{proof}[Proof]
If $I$ is modular, there is $u\in\A$ such that for every
$a\in\A$, we have $a-au\in I$ and $a-ua\in I$. In the
quotient, this is equivalent to
\[
 (a+I)(u+I)=au+I=a+I=ua+I=(u+I)(a+I).
\]
Therefore $u+I$ is a two-sided identity of $\A/I$.
Conversely, if $\A/I$ has an identity, it is a class
$u+I$ for some $u\in\A$. For each $a\in\A$, the
identity laws give $au+I=a+I=ua+I$, that is,
$a-au\in I$ and $a-ua\in I$. This $u$ proves that $I$ is
modular.
\end{proof}

\proofdiagram
{$I\subsetneq\A$ a proper two-sided ideal}
{Definition of multiplication in $\A/I$}
{$u$ modulo $I$ acts as the identity exactly when $a-au,a-ua\in I$}
{$I$ modular $\Longleftrightarrow \A/I$ unital}

\begin{cthejemplo}{(A unital quotient of $C_0(\mathbb R)$)}
To see the proposition in a nonunital algebra, take
$\A=C_0(\mathbb R)$ and
$I_0=\{f\in\A:f(0)=0\}$. The constant function $1$ does not belong
to $\A$, but $u(t)=\max\{1-|t|,0\}$ does and satisfies $u(0)=1$.
Evaluation induces
\[
 \A/I_0\longrightarrow\mathbb C,\qquad f+I_0\longmapsto f(0),
\]
a surjective isomorphism because the functions $\lambda u$ attain
any value $\lambda$. Thus $\A/I_0$ has identity $u+I_0$:
modularity detects an identity that appears only after taking the quotient.
\end{cthejemplo}
